\documentclass[conference]{IEEEtran}
\IEEEoverridecommandlockouts

\usepackage{cite}
\usepackage{amsmath,amssymb,amsfonts}
\usepackage{algorithmic}
\usepackage{graphicx}
\usepackage{textcomp}
\usepackage{xcolor}
\usepackage[utf8]{vietnam}
\usepackage[vietnamese, english]{babel}
\usepackage{booktabs}
\usepackage{tabularx}
\usepackage{array}
\usepackage{multirow}
\def\BibTeX{{\rm B\kern-.05em{\sc i\kern-.025em b}\kern-.08em
    T\kern-.1667em\lower.7ex\hbox{E}\kern-.125emX}}

% Đánh dấu chỗ cần điền (chữ đỏ). Xóa lệnh này khi nộp: thay \todo{...} bằng nội dung thật.
\newcommand{\todo}[1]{\textcolor{red}{[#1]}}

\begin{document}
\selectlanguage{vietnamese}

\title{Cài đặt thuật toán tìm kiếm và tìm kiếm đối kháng\\cho Pac-Man (CS188 Project 1 và Project 2)}

\author{\IEEEauthorblockN{Trường Ngọc Anh}
\IEEEauthorblockA{\textit{Mã sinh viên:} \todo{MSSV}}
\and
\IEEEauthorblockN{Giáp Vũ Diệu Kỳ}
\IEEEauthorblockA{\textit{Mã sinh viên:} \todo{MSSV}}
\and
\IEEEauthorblockN{Nguyễn Mai Phương Thúy}
\IEEEauthorblockA{\textit{Mã sinh viên:} \todo{MSSV}}
\and
\IEEEauthorblockN{Nguyễn Hồng Đăng}
\IEEEauthorblockA{\textit{Mã sinh viên:} 24022956}
}

\maketitle

\renewcommand{\abstractname}{Tóm tắt}
\renewcommand{\figurename}{Hình}
\renewcommand{\tablename}{Bảng}
\renewcommand{\IEEEkeywordsname}{Từ khoá}
\renewcommand{\refname}{Tài liệu tham khảo}

\begin{abstract}
Báo cáo trình bày việc cài đặt các thuật toán tìm kiếm (DFS, BFS, UCS, A*) và tìm kiếm đối kháng (Minimax, Alpha-Beta, Expectimax) để điều khiển Pac-Man trong bộ bài tập CS188 (Project 1 và Project 2). \todo{Thêm 2-3 câu tóm tắt kết quả chính: điểm Autograder, số nút mở rộng, tỉ lệ thắng.}
Mã nguồn: \url{https://github.com/\todo{tai-khoan}/\todo{ten-repo}}.
\end{abstract}

\begin{IEEEkeywords}
Pac-Man, tìm kiếm, A*, heuristic, Minimax, Alpha-Beta, Expectimax.
\end{IEEEkeywords}

\section{Giới thiệu}

\subsection{Tổng quan và mô tả bài toán}
Pac-Man là môi trường mẫu trong môn học CS188 của UC Berkeley \cite{b6}. Báo cáo gồm hai phần:
\begin{itemize}
\item \textbf{Project 1 (Search):} tìm đường đi cho Pac-Man trên mê cung tĩnh, từ tìm đến một điểm cố định đến ăn hết thức ăn, gồm 8 câu (q1--q8).
\item \textbf{Project 2 (Multiagent Search):} Pac-Man đối đầu với các ma, gồm 5 câu (q1--q5).
\end{itemize}
\todo{Nêu rõ vấn đề cần giải quyết: Project 1 là bài toán tìm kiếm trên không gian trạng thái nào, mục tiêu tối ưu gì; Project 2 là bài toán ra quyết định đối kháng với tiêu chí gì. Khoảng 1 đoạn mỗi project.}

\subsection{Phân công công việc}
\begin{table}[!ht]
    \centering
    \renewcommand{\arraystretch}{1.15}
    \scriptsize
    \caption{Bảng phân công công việc cho từng thành viên.}
    \label{tab:phancong}
    \resizebox{\columnwidth}{!}{%
    \begin{tabular}{cccp{4.6cm}}
    \hline
        \textbf{Thành viên} & \textbf{Mã sinh viên} & \textbf{Đóng góp} & \textbf{Công việc} \\
    \hline
        Trường Ngọc Anh & \todo{MSSV} & 25\% & P1: q1--q4 (DFS, BFS, UCS, A*), q8 (ClosestDot); testcase so sánh Minimax và Alpha-Beta \\
        Giáp Vũ Diệu Kỳ & \todo{MSSV} & 25\% & P1: q5 (CornersProblem), q6, q7 (heuristic); P2: q5 (cùng TV4) \\
        Nguyễn Mai Phương Thúy & \todo{MSSV} & 25\% & P2: q1 (Reflex), q2 (Minimax), q3 (Alpha-Beta), q4 (Expectimax) \\
        Nguyễn Hồng Đăng & 24022956 & 25\% & P0: môi trường, repo Git, Autograder; script kiểm thử; P2: q5 (cùng TV2); tổng hợp báo cáo \\
    \hline
    \end{tabular}}
\end{table}

\subsection{Bố cục}
\todo{1-2 câu: Mục II trình bày phương pháp, Mục III trình bày thực nghiệm và kết quả, Mục IV kết luận.}

\section{Phương pháp thực hiện}

\subsection{Project 1: Tìm kiếm}

\subsubsection{Các thuật toán tìm kiếm cơ bản (q1--q4)}
Cả bốn thuật toán cùng dùng khung tìm kiếm đồ thị (graph search) với tập nút đã thăm, chỉ khác cấu trúc lưu biên: \texttt{Stack} cho DFS, \texttt{Queue} cho BFS, \texttt{PriorityQueue} theo chi phí $g(n)$ cho UCS, và theo $f(n)=g(n)+h(n)$ cho A* \cite{b1,b2,b5}. Kiểm tra đích khi lấy nút ra khỏi biên.
\begin{equation}
f(n)=g(n)+h(n)\label{eq:astar}
\end{equation}
\todo{Nêu thêm: cách lưu đường đi (hành động), xử lý nút trùng, tính tối ưu và độ phức tạp của từng thuật toán.}

\subsubsection{Bài toán bốn góc và heuristic (q5, q6)}
\todo{Biểu diễn trạng thái của CornersProblem (vị trí + các góc đã thăm), hàm kế tiếp, kiểm tra đích. Mô tả heuristic cho q6 và chứng minh admissible, consistent.}

\subsubsection{Heuristic cho bài toán ăn hết thức ăn (q7)}
\todo{Mô tả foodHeuristic (ví dụ dùng khoảng cách mê cung \texttt{mazeDistance} giữa các thức ăn xa nhất), chứng minh admissible, consistent; nêu cách lưu đệm trong \texttt{problem.heuristicInfo}.}

\subsubsection{Tìm thức ăn gần nhất (q8)}
\todo{Cách dùng BFS trên \texttt{AnyFoodSearchProblem} để lặp tìm thức ăn gần nhất; nhận xét đường đi tham lam không đảm bảo tối ưu toàn cục.}

\subsection{Project 2: Tìm kiếm đối kháng}

\subsubsection{ReflexAgent (q1)}
\todo{Các đặc trưng trong hàm đánh giá: khoảng cách đến thức ăn gần nhất, khoảng cách đến ma, thời gian ma bị dọa; trọng số.}

\subsubsection{Minimax và Alpha-Beta (q2, q3)}
Minimax luân phiên tầng max (Pac-Man, \texttt{agentIndex}=0) và các tầng min (mỗi con ma một tầng); một mức độ sâu tính trọn một lượt của mọi tác tử \cite{b1,b3}. Alpha-Beta cắt tỉa các nhánh không ảnh hưởng đến giá trị tại gốc, nên cho cùng kết quả với Minimax nhưng sinh ít nút hơn \cite{b4}.
\todo{Công thức/giả mã cho hai thuật toán (có thể dùng gói \texttt{algorithmic}); điều kiện cắt tỉa $\alpha$, $\beta$.}

\subsubsection{Expectimax (q4)}
\todo{Mô hình ma chọn hành động ngẫu nhiên đều; tầng kỳ vọng thay cho tầng min; khác biệt hành vi so với Minimax.}

\subsubsection{Hàm đánh giá tốt hơn (q5)}
\todo{Các đặc trưng đã chọn (thức ăn, capsule, ma thường, ma bị dọa), công thức tổ hợp tuyến tính, cách chọn trọng số.}

\section{Thực nghiệm và kết quả}

\subsection{Thiết lập thực nghiệm}
\todo{Môi trường (Python, hệ điều hành), layout dùng, seed, số ván mỗi cấu hình, tiêu chí đo (số nút mở rộng, độ dài/chi phí đường đi, tỉ lệ thắng, điểm trung bình, thời gian).}

\subsection{Kết quả Project 1}
\begin{table}[!ht]
    \centering
    \caption{Số nút mở rộng và chi phí đường đi trên các mê cung.}
    \label{tab:p1}
    \begin{tabular}{lcccc}
    \toprule
    Layout & DFS & BFS & UCS & A* \\
    \midrule
    \texttt{mediumMaze} & \todo{} & \todo{} & \todo{} & \todo{} \\
    \texttt{bigMaze} & \todo{} & \todo{} & \todo{} & \todo{} \\
    \bottomrule
    \end{tabular}
\end{table}

\begin{table}[!ht]
    \centering
    \caption{Số nút mở rộng của A* với các heuristic (q6, q7).}
    \label{tab:heur}
    \begin{tabular}{lcc}
    \toprule
    Bài toán & Heuristic rỗng & Heuristic của nhóm \\
    \midrule
    Corners (\texttt{mediumCorners}) & \todo{} & \todo{} \\
    Food (\texttt{trickySearch}) & \todo{} & \todo{} \\
    \bottomrule
    \end{tabular}
\end{table}
\todo{Nhận xét ngắn: thuật toán nào mở rộng ít nút nhất, heuristic cải thiện bao nhiêu.}

\subsection{Kết quả Project 2}
\begin{table}[!ht]
    \centering
    \caption{So sánh các tác tử trên \texttt{smallClassic} ($n$ ván).}
    \label{tab:p2}
    \begin{tabular}{lccc}
    \toprule
    Tác tử & Tỉ lệ thắng & Điểm TB & Thời gian/ván \\
    \midrule
    Reflex & \todo{} & \todo{} & \todo{} \\
    Minimax (depth 3) & \todo{} & \todo{} & \todo{} \\
    Alpha-Beta (depth 3) & \todo{} & \todo{} & \todo{} \\
    Expectimax (depth 3) & \todo{} & \todo{} & \todo{} \\
    Hàm đánh giá tốt hơn & \todo{} & \todo{} & \todo{} \\
    \bottomrule
    \end{tabular}
\end{table}
\todo{Nhận xét: Alpha-Beta giảm thời gian bao nhiêu so với Minimax với cùng kết quả; Expectimax so với Minimax khi ma ngẫu nhiên.}

\subsection{Kết quả Autograder}
\begin{table}[!ht]
    \centering
    \caption{Điểm Autograder.}
    \label{tab:grader}
    \begin{tabular}{lcc}
    \toprule
    Phần & Điểm đạt & Điểm tối đa \\
    \midrule
    Project 1 (q1--q8) & \todo{} & 25 \\
    Project 2 (q1--q5) & \todo{} & 25 \\
    \bottomrule
    \end{tabular}
\end{table}
% Có thể thêm hình: \begin{figure}[htbp]\centerline{\includegraphics[width=\columnwidth]{hinh1.png}}\caption{...}\label{fig:hinh1}\end{figure}

\section{Kết luận}
\todo{Kết quả đạt được, hạn chế, hướng phát triển (3-5 câu).}

\begin{thebibliography}{00}
\bibitem{b1} S. Russell and P. Norvig, \textit{Artificial Intelligence: A Modern Approach}, 4th ed. Hoboken, NJ: Pearson, 2021.
\bibitem{b2} P. E. Hart, N. J. Nilsson, and B. Raphael, ``A formal basis for the heuristic determination of minimum cost paths,'' \textit{IEEE Trans. Syst. Sci. Cybern.}, vol. 4, no. 2, pp. 100--107, Jul. 1968.
\bibitem{b3} C. E. Shannon, ``Programming a computer for playing chess,'' \textit{Philos. Mag.}, vol. 41, no. 314, pp. 256--275, 1950.
\bibitem{b4} D. E. Knuth and R. W. Moore, ``An analysis of alpha-beta pruning,'' \textit{Artif. Intell.}, vol. 6, no. 4, pp. 293--326, 1975.
\bibitem{b5} E. W. Dijkstra, ``A note on two problems in connexion with graphs,'' \textit{Numer. Math.}, vol. 1, pp. 269--271, 1959.
\bibitem{b6} J. DeNero and D. Klein, ``Teaching introductory artificial intelligence with Pac-Man,'' in \textit{Proc. 1st Symp. Educ. Adv. Artif. Intell. (EAAI)}, 2010.
\end{thebibliography}

\end{document}